\section{Generalisation to MNLI}\label{sec:mnli}

To test whether the placement and mapping results depend on SST-2, we repeated part of the study on MNLI,
a three-class natural language inference task~\cite{williams2018mnli,wang2019glue}. We used public
BERT-Base\footnote{\nolinkurl{textattack/bert-base-uncased-MNLI}} and
DistilBERT\footnote{\nolinkurl{typeform/distilbert-base-uncased-mnli}} checkpoints with the same
simulator, scores and mapping code. Each model was factorised and fine-tuned for 2 epochs on 100{,}000 MNLI
training pairs, then evaluated on a fixed subset of 2{,}000 matched validation pairs. We ran the
SLC-share study at $\sigma=5\%$ (20 chips), the stuck-cell sweep (20 chips) and fault-aware mapping (10
chips); training, other noise levels and the remaining error sources were not run. \Cref{tab:mnli}
summarises the results.

\begin{table}[t]
\centering\footnotesize
\caption{MNLI accuracy (\%) on 2{,}000 matched validation pairs (chance 33.3\%). Hybrids keep 20\% of the
components in SLC; paired differences are in points over all SLC shares (100 chips).}\label{tab:mnli}
\setlength{\tabcolsep}{4pt}
\begin{tabular}{@{}llcc@{}}
\toprule
& & BERT-Base & DistilBERT \\
\midrule
\multicolumn{2}{@{}l}{Digital SVD model} & 83.2 & 80.9 \\
\multirow{2}{*}{Error-free chip} & All SLC & 72.2 & 76.9 \\
 & Magnitude & 79.1 & 77.3 \\
\midrule
\multirow{4}{*}{$\sigma=5\%$} & Magnitude & \textbf{68.0} & \textbf{71.5} \\
 & Fisher & 67.8 & 71.1 \\
 & Gradient & 63.3 & 71.3 \\
 & Random & 35.9 & 44.1 \\
\addlinespace[2pt]
\multicolumn{2}{@{}l}{\emph{Fisher minus magnitude}} & $-0.2\pm0.2$ & $-0.2\pm0.1$ \\
\multicolumn{2}{@{}l}{\emph{Gradient minus magnitude}} & $-3.4\pm0.4$ & $-0.5\pm0.2$ \\
\midrule
\multirow{3}{*}{0.5\% stuck} & Magnitude & 40.0 & 49.1 \\
 & + fault-aware mapping & \textbf{60.0} & \textbf{67.9} \\
 & + joint mapping & 60.4 & 67.0 \\
\bottomrule
\end{tabular}
\end{table}

\textbf{The placement results hold.} Fisher ranks the components almost as magnitude does (rank
correlation 0.96--0.97, against 0.56--0.57 for the gradient) and reaches the same accuracy on the same
chips, within a fifth of a point. On BERT-Base the gradient is 3.4 points behind; on DistilBERT the
three scores are close, and the gradient is ahead only at 50\% SLC. A random 20\% recovers little.

\textbf{The mapping results hold.} MNLI is more sensitive to stuck cells than SST-2: all-MLC falls to
37--45\% at 0.1\% stuck cells, and the magnitude hybrid to 40--50\% at 0.5\%. Fault-aware mapping recovers much of this,
by $20.1\pm3.0$ points on BERT-Base and $18.8\pm3.2$ on DistilBERT at 0.5\%. Choosing the tier per chip
again adds nothing (within one point of fault-aware mapping).

\textbf{A limitation of the simulator.} On MNLI the error-free chip loses accuracy against the digital
model, by up to 11 points for BERT-Base, and the loss depends on the placement. The cause is the vector
between the two crossbars of a tier, which holds one value per component and is re-quantised to 8 bits
with a single scale. The largest components set that scale, so the outputs of small components are
rounded away. A single tier, as in all-SLC or all-MLC, suffers most (72.2\% and 72.4\% on BERT-Base); a
random split helps little (73.0\%); separating large from small components by magnitude gives each tier
its own scale and recovers most of the loss (79.1\%). On SST-2 the same effect stays within about a point.
Because of it, MNLI comparisons with all-SLC or all-MLC are not meaningful, and we compare only placements
and mappings with the same tier sizes, which share this quantisation.
